\documentclass[10pt]{article}

\usepackage[utf8]{inputenc}

\usepackage[T1]{fontenc}

\usepackage{amsmath}

\usepackage{amsfonts}

\usepackage{amssymb}

\usepackage[version=4]{mhchem}

\usepackage{stmaryrd}

\usepackage{graphicx}

\usepackage[export]{adjustbox}

\graphicspath{ {./images/} }



\begin{document}

чается а т т р а к т о р Л о з и (см. [3]), существование к-рого может быть строго доказано, но в әтом примере гладкость динамич. системы кое-где нарушается.



Лum.: [1] М а р с д ен Д ж., М а к-К р а к ен М., Бифуркаџия рождения цикла и ее приложения, пер. с англ., М.; 1980; [2] Странные аттракторы. Сб. ст., пер. с англ., M., 1981; [3] L o z i R., «J. de Physique», 1978, ser. C, t. 39, № 5, p. 9-10.



СТРАТЕГИЯ в т е о р и и и г р-возможный в соответствии с правилами стратегич. игры способ действия игрока или коалиции действия (см. Иәр теория). В играх в но р м а ль но й ф о р м е (cм. Бескоaлиционная игра) непосредственное описание множеств стратегий входит в «правила» игры. В позиционных играх (см. также Динамическая игра) С. не задаются непосредственно правилами игры, а определяются на их основе косвенным образом. Если в бескоалиционной игре выбор стратегий фиксирован (напр., нек-рым п р инци пом опт има льности), то игра превращается в нестратегическую (см. Kооперативная игра).



СТРАТИФИКАЦИЯ - разложение многообразия (может быть, бесконечномерного) на связные подмногообразия, размерности к-рых строго убывают.



М. И. Войцеховский. СТРЕЛЫБЫ МЕТОД - численный метод решения краевой задачи для обыкновенного дифференциального уравнения, основанный на нек-рой специальной параметризации задачи. См. Пристрелки метод.



СТРОГАЯ ЭРГОДИЧНОСТЬ топологической динамической системь в узком смысле (потока или каскада) - свойство, рассматриваемое в әргодической теоpuu. Оно состоит в следующем: 1) система имеет единственную инвариантную нормированную меру $\mu$ (совместимую с топологией); 2) $\mu(U)>0$ для любого непустого открытого множества $U ; 3)$ для любой ограниченной непрерывной функции $f$ ее временны́ средние вдоль любой траектории стремятся к $\int f d \mu$. Хотя приведенное определение имеет смысл независимо от какихлибо ограничений на фазовое пространство $W$, практически оно применяется тогда, когда $W$ - полное сепарабельное метрич. пространство (обычно даже метрич. компакт). В этом случае С. э. означает, что $W$ представляет собой әргодическое множество. Из С. э. следует, что $W$ является минимальныл множеством (но не обратно). Любой әргодич. поток или каскад в Лебега пространстве метрически изоморфен нек-рому топологич. потоку или каскаду со свойством С. э.



Иногда под С. э. понимается одно только свойство 1); в этом смысле употребляется также термин о д н о ә рг о д и ч но с т ь. Д. В. Аносов.



СТРОГОЙ ИМПЛИКАЦИИ ИСЧИСЛЕНИЕ - логИческое исчисление, основанное на с т р о г о й и мп л и к а ц и и, т. е. логической операции, соответствующей союзу «если..., то...». Для строгой импликации устраняются (полностью или частично) т. н. «парадоксы (материальной) импликации»: ложное высказывание имплицирует (влечет) любое высказывание и истинное высказывание имплицируется любым высказыванием.



С. и. и. ставит своей целью отразить связь по смыслу между посылкой и заключением условного предложения. Существует целый ряд С. и. и. (исчисления Льюиса, Аккермана и т. д.), отличающихся друг от друга тем, что в одних выводимы формулы, невыводимые в других (напр., в исчислениях Льюиса «парадоксы имшликации» устраняются лишь частично, в то время как в исчиєлении Аккермана они устраняются полностью). С. и. и. тесно связано с формализацией модальных выражений («возможно», «невозможно», «необходимо» И т. Д.): в одних исчислениях строгая импликация выражается через модальности, а в других, наоборот, модальности выражаются через строгую импликацию. Лит.: [1] Ф е й с Р., Модальная логика, пер. с англ., M., 1974; [2] L e w i s C. I., L a ng f or d C. H., Symbolic logic 1956, v. 21, №2, p. 113-28.



СТРОФОИДА - плоская алгебраич. кривая 3-го порядка, уравнение к-рой в декартовых прямоугольных координатах $A_{\alpha}$ имеет вид:



\[

y^{2}=x^{2} \frac{d+x}{d-x}

\]



в полярных координатах:



\[

\rho=-d \frac{\cos 2 \varphi}{\cos \varphi}

\]



Начало координат - узловая точка с касательными $y= \pm x$ (см. рис.). Асимптота $x=d$. Площадь петли:



\[

S=2 d^{2}-1 / 2 \pi d^{2} \text {. }

\]



\begin{center}

\includegraphics[max width=\textwidth]{2023_07_09_d50e13efd4ebb1b5c055g-1}

\end{center}



Площадь между кривой и асимптотой:



\[

S_{2}=2 d^{2}+1 / 2 \pi d^{2}

\]



С. относится к т. Н. узлам.



Лли.: [1] С а в е л о в А. А., Плоские кривые, М., 1960;



\includegraphics[max width=\textwidth, center]{2023_07_09_d50e13efd4ebb1b5c055g-1(1)}

вочник по теории кривых третьего порядка, М., 1961 Д.



СТРУВЕ ФУНКЦИЯ - Функция



\[

\begin{gathered}

\mathbf{H}_{v}(z)=\frac{2}{\sqrt{\pi}} \frac{(z / 2)^{v}}{\Gamma(v+1 / 2)} \int_{0}^{\pi / 2} \sin (z \cos t) \sin ^{2} v t d t, \\

\operatorname{Re} v>-1 / 2,

\end{gathered}

\]



удовлетворяющая неоднородному уравнению Бесселя:



\[

z^{2} y^{\prime \prime}+z y^{\prime}+\left(z^{2}-v^{2}\right)_{y}=\frac{4(z / 2)^{v+1}}{\sqrt{\pi} \Gamma(v+1 / 2)} .

\]



Разложение в степенной ряд:



\[

\begin{gathered}

\mathbf{H}_{v}(z)= \\

=\frac{2}{\sqrt{\pi}} \frac{(z / 2)^{v+1}}{\Gamma(v+3 / 2)}\left[1-\frac{z^{2}}{3(2 v+3)}+\frac{z^{4}}{3 \cdot 5(2 v+3)(2 v+5)}-\ldots\right] .

\end{gathered}

\]



С. Ф. целого порядка $n$ связана с Вебера фуункцией следующими соотношениями:



\[

\begin{gathered}

\mathbf{H}_{0}(z)=-E_{0}(z), \mathbf{H}_{1}(z)=2 / \pi-E_{1}(z), \\

\mathbf{H}_{2}(z)=\frac{z}{3 \pi}-E_{2}(z) \ldots .

\end{gathered}

\]



С. ф. порядка $n+1 / 2(n-$ целое число) являются әлементарными функциями, напр.



\[

\sqrt{\frac{\pi}{\alpha}} \mathbf{H}_{1 / 2}(z)=\frac{1-\cos z}{\sqrt{z}} .

\]



При $|z| \gg 1,|z| \gg|v|$ имеет место асимптотич. разложение



\[

\begin{gathered}

\mathbf{H}_{v}(z)-N_{v}(z) \approx \\

\approx \frac{(z / 2)^{v-1}}{\sqrt{\pi} \Gamma(v+1 / 2)}\left[1+\frac{1 \cdot(2 v-1)}{z^{2}}+\frac{1 \cdot 3 \cdot(2 v-1)(2 v-3)}{z^{4}}+\ldots\right],

\end{gathered}

\]



где $N_{v}(z)$-Неймана функция.



М о дифици р о в а нн а я



ф $\mathbf{н}$ к ц я С т р у в е - функция



\[

\mathbf{L}_{v}(z)=-i e^{i v \pi / 2} \mathbf{H}_{v}(i z) .

\]



Ее разложение в ряд:



\[

\mathbf{L}_{v}(z)=\left(\frac{z}{2}\right)^{v+1} \sum_{k=0}^{\infty} \frac{(z / 2)^{2 k}}{\Gamma(k+3 / 2) \Gamma(k+v+3 / 2)} .

\]



Для больших $|z|$ имеет место асимптотич. разложение



\[

\mathbf{L}_{v}(z)-I_{-v}(z) \approx \frac{1}{\pi},

\]



где $I_{-v}$ - модифицированная функция Бесселя.



С. ф. иногда обозначаөтся $S_{v}(z)$. Функцию ввел Х. Струве [1].





\end{document}